\documentclass{report}
\usepackage{luatexja} % LuaTeXで日本語を使うためのパッケージ
\usepackage{luatexja-fontspec} % LuaTeX用の日本語フォント設定

% --- 数学関連 ---
\usepackage{amsmath, amssymb, amsfonts, mathtools, bm, amsthm} % 基本的な数学パッケージ
\usepackage{type1cm, upgreek} % 数式フォントとギリシャ文字
\usepackage{physics, mhchem} % 物理や化学の記号や式の表記を簡単にする

% --- 表関連 ---
\usepackage{multirow, longtable, tabularx, array, colortbl, dcolumn, diagbox} % 表のレイアウトを柔軟にする
\usepackage{tablefootnote, truthtable} % 表中に注釈を追加、真理値表
\usepackage{tabularray} % 高度な表組みレイアウト

% --- グラフィック関連 ---
\usepackage{tikz, graphicx} % 図の描画と画像の挿入
\usepackage{background} % ウォーターマークの設定
\usepackage{caption, subcaption} % 図や表のキャプション設定
\usepackage{float, here} % 図や表の位置指定

% --- レイアウトとページ設定 ---
\usepackage{fancyhdr} % ページヘッダー、フッター、余白の設定
\usepackage[top = 20truemm, bottom = 20truemm, left = 20truemm, right = 20truemm]{geometry}
\usepackage{fancybox, ascmac} % ボックスのデザイン

% --- 色とスタイル ---
\usepackage{xcolor, color, colortbl, tcolorbox} % 色とカラーボックス
\usepackage{listings, jvlisting} % コードの色付けとフォーマット

% --- 参考文献関連 ---
\usepackage{biblatex, usebib} % 参考文献の管理と挿入
\usepackage{url, hyperref} % URLとリンクの設定

% --- その他の便利なパッケージ ---
\usepackage{footmisc} % 脚注のカスタマイズ
\usepackage{multicol} % 複数段組
\usepackage{comment} % コメントアウトの拡張
\usepackage{siunitx} % 単位の表記
\usepackage{docmute}
% \usepackage{appendix}
% --- tcolorboxとtikzの設定 ---
\tcbuselibrary{theorems, breakable} % 定理のボックスと改ページ設定
\usetikzlibrary{decorations.markings, arrows.meta, calc} % tikzの装飾や矢印の設定

% --- 定理スタイルと数式設定 ---
\theoremstyle{definition} % 定義スタイル
\numberwithin{equation}{section} % 式番号をサブセクション単位でリセット

% --- hyperrefの設定 ---
\hypersetup{
  setpagesize = false,
  bookmarks = true,
  bookmarksdepth = tocdepth,
  bookmarksnumbered = true,
  colorlinks = false,
  pdftitle = {}, % PDFタイトル
  pdfsubject = {}, % PDFサブジェクト
  pdfauthor = {}, % PDF作者
  pdfkeywords = {} % PDFキーワード
}

% --- siunitxの設定 ---
\sisetup{
  table-format = 1.5, % 小数点以下の桁数
  table-number-alignment = center, % 数値の中央揃え
}

% --- 透かし画像の設定 --- 
\backgroundsetup{
  scale=0.5,                       % 画像のスケール
  % color=black,                   % 画像の色（透かし用に半透明が推奨）
  opacity=0.2,                   % 透かしの透明度（0が完全透明、1が完全不透明）
  angle=0,                       % 画像の角度
  position = current page.south east,  % ページの右下
  hshift=-6cm, % 右方向へのシフト（負の値で内側に移動）
  vshift=5cm,  % 上方向へのシフト（正の値で内側に移動）
  contents={\includegraphics{./fig/appilogo-circular-full.png}} % 画像のパス
}

% --- その他の設定 ---
\allowdisplaybreaks % 数式の途中改ページ許可
\newcolumntype{t}{!{\vrule width 0.1pt}} % 新しいカラムタイプ
\newcolumntype{b}{!{\vrule width 1.5pt}} % 太いカラム
\UseTblrLibrary{amsmath, booktabs, counter, diagbox, functional, hook, html, nameref, siunitx, varwidth, zref} % tabularrayのライブラリ
\setlength{\columnseprule}{0.4pt} % カラム区切り線の太さ
\captionsetup[figure]{font = bf} % 図のキャプションの太字設定
\captionsetup[table]{font = bf} % 表のキャプションの太字設定
\captionsetup[lstlisting]{font = bf} % コードのキャプションの太字設定
\captionsetup[subfigure]{font = bf, labelformat = simple} % サブ図のキャプション設定
\setcounter{secnumdepth}{4} % セクションの深さ設定
\newcolumntype{d}{D{.}{.}{5}} % 数値のカラム
\newcolumntype{M}[1]{>{\centering\arraybackslash}m{#1}} % センター揃えのカラム
\DeclareMathOperator{\diag}{diag}
\everymath{\displaystyle} % 数式のスタイル
\newcommand{\inner}[2]{\left\langle #1, #2 \right\rangle}
\renewcommand{\figurename}{図}
\renewcommand{\i}{\mathrm{i}} % 複素数単位i
\renewcommand{\laplacian}{\grad^2} % ラプラシアンの記号
\renewcommand{\thesubfigure}{(\alph{subfigure})} % サブ図の番号形式
\newcommand{\m}[3]{\multicolumn{#1}{#2}{#3}} % マルチカラムのショートカット
\renewcommand{\r}[1]{\mathrm{#1}} % mathrmのショートカット
\newcommand{\e}{\mathrm{e}} % 自然対数の底e
\newcommand{\Ef}{E_{\mathrm{F}}} % フェルミエネルギー
\renewcommand{\c}{\si{\degreeCelsius}} % 摂氏記号
\renewcommand{\d}{\r{d}} % d記号
\renewcommand{\t}[1]{\texttt{#1}} % タイプライタフォント
\newcommand{\kb}{k_{\mathrm{B}}} % ボルツマン定数
% \renewcommand{\phi}{\varphi} % ϕをφに変更
\renewcommand{\epsilon}{\varepsilon}
\newcommand{\fullref}[1]{\textbf{\ref{#1} \nameref{#1}}}
\newcommand{\reff}[1]{\textbf{図\ref{#1}}} % 図参照のショートカット
\newcommand{\reft}[1]{\textbf{表\ref{#1}}} % 表参照のショートカット
\newcommand{\refe}[1]{\textbf{式\eqref{#1}}} % 式参照のショートカット
\newcommand{\refp}[1]{\textbf{コード\ref{#1}}} % コード参照のショートカット
\renewcommand{\lstlistingname}{コード} % コードリストの名前
\renewcommand{\theequation}{\thesection.\arabic{equation}} % 式番号の形式
\renewcommand{\footrulewidth}{0.4pt} % フッターの線
\newcommand{\mar}[1]{\textcircled{\scriptsize #1}} % 丸囲み文字
\newcommand{\combination}[2]{{}_{#1} \mathrm{C}_{#2}} % 組み合わせ
\newcommand{\thline}{\noalign{\hrule height 0.1pt}} % 細い横線
\newcommand{\bhline}{\noalign{\hrule height 1.5pt}} % 太い横線

% --- カスタム色定義 ---
\definecolor{burgundy}{rgb}{0.5, 0.0, 0.13} % バーガンディ色
\definecolor{charcoal}{rgb}{0.21, 0.27, 0.31} % チャコール色
\definecolor{forest}{rgb}{0.0, 0.35, 0} % 森の緑色

% --- カスタム定理環境の定義 ---
\newtcbtheorem[number within = chapter]{myexc}{練習問題}{
  fonttitle = \gtfamily\sffamily\bfseries\upshape,
  colframe = forest,
  colback = forest!2!white,
  rightrule = 1pt,
  leftrule = 1pt,
  bottomrule = 2pt,
  colbacktitle = forest,
  theorem style = standard,
  breakable,
  arc = 0pt,
}{exc-ref}
\newtcbtheorem[number within = chapter]{myprop}{命題}{
  fonttitle = \gtfamily\sffamily\bfseries\upshape,
  colframe = blue!50!black,
  colback = blue!50!black!2!white,
  rightrule = 1pt,
  leftrule = 1pt,
  bottomrule = 2pt,
  colbacktitle = blue!50!black,
  theorem style = standard,
  breakable,
  arc = 0pt
}{proposition-ref}
\newtcbtheorem[number within = chapter]{myrem}{注意}{
  fonttitle = \gtfamily\sffamily\bfseries\upshape,
  colframe = yellow!20!black,
  colback = yellow!50,
  rightrule = 1pt,
  leftrule = 1pt,
  bottomrule = 2pt,
  colbacktitle = yellow!20!black,
  theorem style = standard,
  breakable,
  arc = 0pt
}{remark-ref}
\newtcbtheorem[number within = chapter]{myex}{例題}{
  fonttitle = \gtfamily\sffamily\bfseries\upshape,
  colframe = black,
  colback = white,
  rightrule = 1pt,
  leftrule = 1pt,
  bottomrule = 2pt,
  colbacktitle = black,
  theorem style = standard,
  breakable,
  arc = 0pt
}{example-ref}
\newtcbtheorem[number within = chapter]{exc}{Requirement}{myexc}{exc-ref}
\newcommand{\rqref}[1]{{\bfseries\sffamily 練習問題 \ref{exc-ref:#1}}}
\newtcbtheorem[number within = chapter]{definition}{Definition}{mydef}{definition-ref}
\newcommand{\dfref}[1]{{\bfseries\sffamily 定義 \ref{definition-ref:#1}}}
\newtcbtheorem[number within = chapter]{prop}{命題}{myprop}{proposition-ref}
\newcommand{\prref}[1]{{\bfseries\sffamily 命題 \ref{proposition-ref:#1}}}
\newtcbtheorem[number within = chapter]{rem}{注意}{myrem}{remark-ref}
\newcommand{\rmref}[1]{{\bfseries\sffamily 注意 \ref{remark-ref:#1}}}
\newtcbtheorem[number within = chapter]{ex}{例題}{myex}{example-ref}
\newcommand{\exref}[1]{{\bfseries\sffamily 例題 \ref{example-ref:#1}}}
% --- 再定義コマンド ---
% \mathtoolsset{showonlyrefs=true} % 必要な式番号のみ表示
\pagestyle{fancy} % ヘッダー・フッターのスタイル設定
\chead{応用量子物性講義ノート} % 中央ヘッダー
% \rhead{}
\fancyhead[R]{\rightmark}
\renewcommand{\sectionmark}[1]{\markright{\thesection\ #1}}
\cfoot{\thepage} % 中央フッターにページ番号
\lhead{}
\rfoot{Yuto Masuda and Haruki Aoki} % 右フッターに名前
\setcounter{tocdepth}{4} % 目次の深さ
\makeatletter
\@addtoreset{equation}{section} % サブセクションごとに式番号をリセット
\makeatother

% --- メタ情報 ---
\title{応用量子物性講義ノート}
\date{更新日\today}
\author{Yuto Masuda and Haruki Aoki}

\begin{document}
  本節では，次節にて行う部分波展開について，数学的な準備を行う．
  最初に，以下で扱う演算子のクラスであるStrum-Liouvillle演算子についての性質を調べる．
  まず，Legendre多項式とLegendre陪多項式の性質を調べて，直交性を知る．
  次に，Bessel微分方程式の解であるBessel関数の表式を求めて，球Bessel関数と球Neumann関数，球Hankel関数を定義する．
  続いて，Legendre多項式を用いて波動関数を展開して，展開係数$B_{ml}$が満たすべき微分方程式を導く．
  最後に，展開係数$B_{ml}$が球Bessel関数と球Neumann関数の線型結合で書けることを確かめて，線型結合の係数が波動関数を特徴づけるものだと知る．
  \subsection{Strum-Liouvillle演算子のHermite性}
    $a < b$として，$x\in \qty[a, b]$で定義された関数空間$V$を考える．$\rho(x)$を非負の実数関数として，$f, g\in V$に対して，
    \begin{align}
      \inner{f}{g} \coloneqq \int_{a}^{b} f^*(x)g(x)\rho(x)\dd{x}
    \end{align}
    なる内積を入れる．
    関数空間$V$上の演算子として$\mathcal{L}$を，
    \begin{align}
      \mathcal{L} \coloneqq \frac{1}{\rho(x)}\qty[\dv{x}\qty{p(x)\dv{x}} + q(x)]\label{strum-liouville-operator-def}
    \end{align}
    とする．
    \refe{strum-liouville-operator-def}なる形をした演算子をStrum-Liouvillle演算子という．
    境界条件を，$\forall f \in V$について，
    \begin{align}
      \begin{dcases}
        f(a) = f(b) \\ 
        p(a)f'(a) = p(b)f'(b)
      \end{dcases}
    \end{align}
    とすると，
    \begin{align}
      \inner{f}{\mathcal{L}g} = \inner{\mathcal{L}f}{g}\label{hermite-def}
    \end{align}
    が成立する．
    \refe{hermite-def}なる関係が成り立つ演算子$\mathcal{L}$をHermite演算子という．
    \begin{proof}
      内積の定義より，
      \begin{align}
        \inner{f}{\mathcal{L}g} &= \int_{a}^{b} f^*(x)\frac{1}{\rho(x)}\qty[\dv{x}\qty{p(x)\dv{x}g(x)} + q(x)g(x)]\rho(x)\dd{x} \\ 
        &= \int_{a}^{b} f^*(x)\qty[\dv{x}\qty{p(x)\dv{x}g(x)} + q(x)g(x)]\dd{x} \\ 
        &= \qty[f^*(x)g'(x)]_{a}^{b} - \int_{a}^{b}f'^*(x)p(x)g'(x)\dd{x} + \int_{a}^{b}f^*(x)q(x)g(x)\dd{x} \\ 
        &= \qty[f^*(x)p(x)g'(x)]_{a}^{b} - \qty[f'^*(x)p(x)g(x)]_{a}^{b} + \int_{a}^{b}g(x)\frac{1}{\rho(x)}\qty[\dv{x}\qty(p(x)\dv{x}f^*(x)) + q(x)f^*(x)]\rho(x)\dd{x} \\ 
        &= \qty[f^*(x)p(x)g'(x)]_{a}^{b} - \qty[f'^*(x)p(x)g(x)]_{a}^{b} + \inner{g}{\mathcal{L}f}^* \\ 
        &= \qty[f^*(x)p(x)g'(x)]_{a}^{b} - \qty[g(x)p(x)f'^*(x)]_{a}^{b} + \inner{\mathcal{L}f}{g} 
      \end{align}
      となる．第1項と第2項について，第1項に$f(a) = f(b)$を，第2項に$p(x)$が実数値関数であり$p(a)f'(a) = p(b)f'(b)$であることを用いると
      \begin{align}
        \qty[f^*(x)p(x)g'(x)]_{a}^{b} - \qty[f'^*(x)p(x)g(x)]_{a}^{b} = \qty{p(b)g'(b) - p(a)f(a)}f^*(a) - \qty{g(b) - g(a)}p(a)f'(a)
      \end{align}
      を得る．今度は，第1項に$p(x)$が実数値関数であり$p(a)g'(a) = p(b)g'(b)$であることを，第2項に$g(a) = g(b)$を用いれば，
      \begin{align}
        \inner{f}{\mathcal{L}g} = \inner{\mathcal{L}f}{g}
      \end{align}
      を得る．
    \end{proof}
  \subsection{Legendre多項式}
    \refe{strum-liouville-operator-def}において，$a = -b = 1$とする．また，
    \begin{align}
      \rho(x) &\coloneqq 1 \\ 
      p(x) &\coloneqq 1 - x^2 \\ 
      q(x) &\coloneqq -\frac{m^2}{1 - x^2},\ m\in \qty{0, 1, \cdots}
    \end{align}
    とすると，
    \begin{align}
      \mathcal{L}_m = \dv{x}\qty{(1 - x^2)\dv{x}} - \frac{m^2}{1 - x^2}
    \end{align}
    となる．
    関数の内積は，
    \begin{align}
      \inner{f}{g} = \int_{-1}^{1} f^*(x)g(x) \dd{x} \label{inner-prod-legendre}
    \end{align}
    と定義しておく．
    また，境界条件は，
    \begin{align}
      p(\pm 1)f^*(\pm 1)g'(\pm 1) = 0
    \end{align}
    とする．これより，演算子$\mathcal{L}$がHermite演算子であると確かめられる．
    たとえば，$g'(x) = p(x)$となるように$g(x)$を定めれば，$f(1) = f(-1) = 0$となる．
    また，$f(x) = 1$となるように$f(x)$を定めれば，$p(1)g'(1) = p(-1)g'(1) = 0$となるので，前節で示した境界条件を満足する．
    さて，Legendre多項式$P_n(x)$は$n$を非負整数として，
    \begin{align}
      \mathcal{L}_0P_n(x) &= -n(n + 1) P_n(x) \\ 
      \Leftrightarrow \dv{x}\qty{(1 - x^2)\dv{x}P_n(x)} &= -n(n + 1)P_n(x)\label{legendre-polynominal} 
    \end{align}
    なる$P_n(x)$のうち，$x = 0$周りで級数展開したもので，
    \begin{align}
      P_n(x) = \sum_{j = 0}^{\infty}u_jx^j\label{legendre-polynominal-u-expantion}
    \end{align}
    と書いたとき，
    \begin{align}
      u_n &= \frac{(2n)!}{2^n(n!)^2}\label{legendre-polynominal-u-expantion-n} \\ 
      u_{n + 1} &= 0\label{legendre-polynominal-u-expantion-np1}
    \end{align}
    なるものである．\refe{legendre-polynominal-u-expantion}を\refe{legendre-polynominal}に代入すると，
    \begin{align}
      \dv{x}\qty{(1 - x^2)\dv{x}\sum_{j = 0}^{\infty}u_jx^j} &= -n(n + 1)\sum_{j = 0}^{\infty}u_jx^j \\ 
      \Leftrightarrow \sum_{j = 0}^{\infty}ju_j\dv{x}\qty(x^{j - 1} - x^{j + 1}) &= -n(n + 1)\sum_{j = 0}^{\infty}u_jx^j \\ 
      \Leftrightarrow \sum_{j = 0}^{\infty}j(j - 1)u_jx^{j - 2} - \sum_{j = 0}^{\infty}j(j + 1)u_jx^j &= -n(n + 1)\sum_{j = 0}^{\infty}u_jx^j \\ 
      \Leftrightarrow \sum_{j = 0}^{\infty}j(j - 1)u_jx^{j - 2} &= \sum_{j = 0}^{\infty}u_j\qty[j(j + 1) - n(n + 1)]x^j \\ 
      \Leftrightarrow \sum_{j = 2}^{\infty}j(j - 1)u_jx^{j - 2} &= \sum_{j = 0}^{\infty}u_j\qty[j(j + 1) - n(n + 1)]x^j \\ 
      \Leftrightarrow \sum_{j = 0}^{\infty}(j + 1)(j + 2)u_{j + 2}x^{j} &= \sum_{j = 0}^{\infty}u_j\qty[j(j + 1) - n(n + 1)]x^j 
    \end{align}
    となるから，
    \begin{align}
      (j + 1)(j + 2)u_{j + 2} = \qty[j(j + 1) - n(n + 1)]u_j\label{legendre-polynominal-u-expantion-recurr}
    \end{align}
    なる漸化式が成立する．
    \refe{legendre-polynominal-u-expantion-recurr}において$j = n$を代入すると，$u_{n + 2} = 0$となる．
    また，$j = n + 1$を代入すると\refe{legendre-polynominal-u-expantion-np1}より$u_{n + 1} = 0$である．
    よって，
    \begin{align}
      0 = u_{n + 1} = u_{n + 2} = u_{n + 3} = u_{n + 4} = \cdots 
    \end{align}
    となる．また，\refe{legendre-polynominal-u-expantion-recurr}に$j = n - 2$を代入すると，
    \begin{align}
      u_{n - 2} = -\frac{n(n - 1)}{2(2n - 1)}u_n\label{legendre-polynominal-u-expantion-recurr-2}
    \end{align}
    となる．よって，\refe{legendre-polynominal-u-expantion-n}と\refe{legendre-polynominal-u-expantion-recurr-2}を用いて
    \refe{legendre-polynominal-u-expantion}を表すと，
    \begin{align}
      P_n(x) &= \frac{(2n)!}{2^n(n!)^2}\qty[x^n - \frac{n(n - 1)}{2(2n - 1)}x^{n - 2} + \frac{n(n - 1)(n - 2)(n - 3)}{2\cdot 4\cdot (2n - 1)(2n - 3)}x^{n - 4} + \cdots] \\ 
      &= \sum_{s = 0}^{\lfloor n/2 \rfloor}(-1)^s \frac{(2n - 2s)!}{2^ns!(n - s)!(n - 2s)!}x^{n - 2s}
    \end{align}
    となる．
    なお，Legendre多項式は，
    \begin{align}
      \dv[n]{x}\qty(x^2 - 1)^n &= \sum_{s = 0}^{n}\dv[n]{x}(-1)^s\mqty(n \\ k)x^{2n - 2k} \\ 
      &= \sum_{s = 0}^{\lfloor n/2 \rfloor}(-1)^s\frac{n!}{s!(n - s)!}\frac{(2n - 2s)!}{(n - 2k)!}
    \end{align}
    なる関係を用いると，
    \begin{align}
      P_n(x) &= \frac{1}{2^nn!}\dv[n]{x}\qty(x^2 - 1)^n\label{legendre-pn-easy}
    \end{align}
    となる．
  \subsection{Legendre陪多項式}
    Legendre陪多項式$P_n^m(x)$は$m \leq n$として，
    \begin{align}
      P_n^m(x) &\coloneqq \qty(1 - x^2)^{m/2}\dv[m]{x}P_n(x)
    \end{align}
    と定義されて，\refe{legendre-pn-easy}を用いれば，
    \begin{align}
      P_n^m(x) &= \qty(1 - x^2)^{m/2}\dv[m]{x}P_n(x) \\ 
      &= \qty(1 - x^2)^{m/2}\dv[m]{x}\frac{1}{2^nn!}\dv[n]{x}\qty(x^2 - 1)^n \\ 
      &= \frac{1}{2^nn!}\qty(1 - x^2)^{m/2}\dv[n + m]{x}\qty(x^2 - 1)^n \label{assos-legendre-polynominal}
    \end{align}
    となる．
    また，$P_n^m$はLegendreの陪微分方程式，
    \begin{align}
      \mathcal{L}_mP_n^m(x) + n(n + 1)P_n^m(x) = 0\label{assos-legendre-diff-eq}
    \end{align}
    を満たす．
    Legendre陪多項式の直交性は$\mathcal{L}_m$がHermite演算子であり，その固有関数である$P_n^m(x)$が直交することより従う．
    自分自身との内積，つまり，$\inner{P_n^m(x)}{P_n^m(x)}$の値を計算する．
    \refe{assos-legendre-polynominal}を用いて，\refe{inner-prod-legendre}で示した内積の定義に従って計算する．
    $n + m$回部分積分を行うと，
    \begin{align}
      \inner{P_{n}^m(x)}{P_{n}^m(x)} &= \int_{-1}^{1}\frac{1}{2^{2n}\qty(n!)^2}\qty(1 - x^2)^m \qty{\dv[n + m]{x}\qty(x^2 - 1)^m}\qty{\dv[n + m]{x}\qty(x^2 - 1)^m} \dd{x} \\ 
      &= \frac{1}{2^{2n}\qty(n!)^2}\int_{-1}^{1}\qty(1 - x^2)^m \qty{\dv[n + m]{x}\qty(x^2 - 1)^m}\dv{x}\qty{\dv[n + m - 1]{x}\qty(x^2 - 1)^m} \dd{x} \\ 
      &= \frac{1}{2^{2n}\qty(n!)^2}\qty[\qty(1 - x^2)^m\qty{\dv[n + m]{x}\qty(x^2 - 1)^m}\qty{\dv[n + m - 1]{x}\qty(x^2 - 1)^m}]_{-1}^{1} \\ 
      &\ \ -\frac{1}{2^{2n}\qty(n!)^2} \int_{-1}^{1}\qty{\dv{x}\qty(1 - x^2)^m\qty{\dv[n + m]{x}(x^2 - 1)^n}}\qty{\dv[n + m + 1]{x}\qty(x^2 - 1)^m}\dd{x} \\ 
      &= -\frac{1}{2^{2n}\qty(n!)^2}\int_{-1}^{1}\qty{\dv{x}\qty(1 - x^2)^m\qty{\dv[n + m]{x}(x^2 - 1)^n}}\qty{\dv[n + m + 1]{x}\qty(x^2 - 1)^m} \dd{x} \\
      &= \cdots \\ 
      &= \frac{\qty(-1)^{n + m}}{2^{2n}\qty(n!)^2}\int_{-1}^{1}\qty(x^2 - 1)^n\qty{\dv[n + m]{x}\qty(1 - x^2)^m\qty{\dv[n + m]{x}\qty(x^2 - 1)^n}} \dd{x} \\ 
      &= \frac{\qty(-1)^{n + m}}{2^{2n}\qty(n!)^2}\int_{-1}^{1}\qty(x^2 - 1)^n\qty[\sum_{k = 0}^{n + m}\mqty(n + m \\ k)\qty{\dv[n + m - k]{x}\qty(1 - x^2)^m}\qty{\dv[n + m + k]{x}\qty(x^2 - 1)^n}] \dd{x}\label{use-leibniz-formula}
    \end{align}
    となる．最終行でLeibnizの公式を用いた．
    \refe{use-leibniz-formula}の和の中の$n + m - k$階微分と$n + m - k$階微分を考える．
    $\qty(1 - x^2)^m$と$(x^2 - 1)^n$の最高次数は，それぞれ$2m$と$2n$であるから，$2m \geq n + m - k$かつ$2n \geq n + m + k$なる$k$でのみ和の中は0でなくなる．
    つまり，$n - m\leq k$かつ$n - m \geq k$なる$k$は$k = n - m$のみである．
    よって，\refe{use-leibniz-formula}は，
    \begin{align}
      \inner{P_{n}^m(x)}{P_{n}^m(x)} &= \frac{\qty(-1)^{n + m}}{2^{2n}\qty(n!)^2}\int_{-1}^{1}\qty(x^2 - 1)^n\qty[\sum_{k = 0}^{n + m}\mqty(n + m \\ k)\qty{\dv[n + m - k]{x}\qty(1 - x^2)^m}\qty{\dv[n + m + k]{x}\qty(x^2 - 1)^n}] \dd{x} \\ 
      &= \frac{\qty(-1)^{n + m}}{2^{2n}\qty(n!)^2}\int_{-1}^{1}\qty(x^2 - 1)^n\qty[\mqty(n + m \\ n - m)\qty{\dv[2m]{x}\qty(1 - x^2)^m}\qty{\dv[2n]{x}\qty(x^2 - 1)^n}] \dd{x} \\ 
      &= \frac{\qty(-1)^{n + m}}{2^{2n}\qty(n!)^2}(-1)^m\qty(2m)!\qty(2n)!\frac{(n + m)!}{(n - m)!(2m)!}\int_{-1}^{1}\qty(x^2 - 1)^n\dd{x} 
    \end{align}
    積分は，$x = \cos\theta$と置換すると，
    \begin{align}
      \int_{-1}^{1}\qty(x^2 - 1)^n\dd{x} = (-1)n\int_{-1}^{1}\sin^{2n + 1}\theta\dd{\theta}
    \end{align}
    となる．$I_{2n + 1}$を，
    \begin{align}
      I_{2n + 1} \coloneqq \int_{0}^{\pi}\sin^{2n + 1}\theta\dd{\theta}
    \end{align}
    と定義すると，
    \begin{align}
      I_{2n + 1} &= \int_{0}^{\pi} \sin^{2n}\theta\dv{\theta}(-\cos\theta)\dd{\theta} \\ 
      &= \qty[\sin^{2n}\theta\cdot(-\cos\theta)]_{0}^{\pi} + 2n\int_{0}^{\pi}\sin^{2n - 1}\theta \cos^2\theta\dd{\theta} \\ 
      &= 2nI_{2n - 1} - 2nI_{2n + 1}
    \end{align}
    となるので，
    \begin{align}
      I_{2n + 1} &= \frac{2n}{2n + 1}I_{2n - 1} \\ 
      &= \frac{2n}{2n + 1}\frac{2n - 2}{2n - 1}I_{2n - 3} \\ 
      &= \cdots \\ 
      &= \frac{(2n)!!}{(2n + 1)!!}I_1 \\ 
      &= 2\frac{(2n)!!}{(2n + 1)!!} \\ 
      &= 2\cdot 2^nn! \cdot \frac{2^nn!}{(2n + 1)!}
    \end{align}
    となる．2重階乗について，
    \begin{align}
      (2n)!! &= 2^nn! \\ 
      (2n + 1) !! &= \frac{(2n + 1)!}{(2n)!!} = \frac{(2n + 1)!}{2^nn!}
    \end{align}
    なる関係が成り立つことを用いると，
    \begin{align}
      \inner{P_{n}^m(x)}{P_{n}^m(x)} &= \frac{\qty(-1)^{n + m}}{2^{2n}\qty(n!)^2}(-1)^m\qty(2m)!\qty(2n)!\frac{(n + m)!}{(n - m)!(2m)!}(-1)^n2\frac{(2n)!!}{(2n + 1)!!} \\
      &=  \frac{\qty(-1)^{n + m}}{2^{2n}\qty(n!)^2}(-1)^m\qty(2m)!\qty(2n)!\frac{(n + m)!}{(n - m)!(2m)!}(-1)^n2\cdot 2^nn! \cdot \frac{2^nn!}{(2n + 1)!} \\
      &= \frac{2}{2n + 1}\frac{(n + m)!}{(n - m)!}
    \end{align}
    となる．
    Legendre多項式の直交性とまとめて書くと，
    \begin{align}
      \inner{P_{n'}^m(x)}{P_{n}^m(x)} = \delta_{n}^{n'}\frac{2}{2n + 1}\frac{(n + m)!}{(n - m)!}\label{pnm-norm}
    \end{align}
    となる．
  \subsection{Bessel関数}
    Legendre陪多項式を考えていたときとは別のStrum-Liouvillle演算子を考える．
    \refe{strum-liouville-operator-def}において，$a \to 0$，$b \to a$とする．また，
    \begin{align}
      \rho(z) &\coloneqq x \\ 
      p(z) &\coloneqq x \\ 
      q(z) &\coloneqq -\frac{\nu^2}{z},\ \nu \geq 0
    \end{align}
    とする．
    すなわち，
    \begin{align}
      \mathcal{L}_{\nu} &= \frac{1}{z}\qty[\dv{z}\qty{z\dv{z}} - \frac{\nu^2}{z}] \\ 
      \inner{f}{g} &= \int_{0}^{a}f^*(z)g(z)x\dd{z}
    \end{align}
    である．
    さて，
    \begin{align}
      \mathcal{L}_{\nu}J_{\nu}(z) &= - J_{\nu}(z)  \\ 
      \Leftrightarrow \qty(\dv[2]{z} + \frac{1}{z}\dv{z} + 1 - \frac{\nu^2}{z^2}) &= 0 \label{bessel-differential-eq-sub}\\ 
      \Leftrightarrow \dv[2]{z}J_{\nu}(z) + \frac{1}{z}\dv{z}J_{\nu}(z) - \frac{\nu^2}{z^2}J_{\nu}(z) &= J_{\nu}(z) \label{bessel-differential-eq}
    \end{align}
    なる$J_{\nu}(z)$を考える．1階微分の項と微分をしない項は$z = 0$で発散するが，それぞれ$z$と$z^2$をかければ発散しないので確定特異点である．
    このとき，$z = 0$の周りで$J_{\nu}(z)$を指数$\alpha$のFrobenius展開をすると，
    \begin{align}
      J_{\nu}(z) = z^{\alpha}\sum_{n}u_nz^n\label{bessel-func-frobenius}
    \end{align}
    となる．
    ただし，計算の簡単のために$n < 0$なる任意の$n$に対して$u_n = 0$と定めた．
    また，$u_0 \neq 0$とする．
    \refe{bessel-func-frobenius}を\refe{bessel-differential-eq}に代入することを考える．
    $J_{\nu}(z)$の1階微分と2階微分が，
    \begin{align}
      \dv{z}J_{\nu}(z) &= \sum_{n}(n + \alpha)u_nz^{n + \alpha - 1} \\ 
      \dv[2]{z}J_{\nu}(z) &= \sum_{n}(n + \alpha)(n + \alpha - 1)u_nz^{n + \alpha - 2} 
    \end{align}
    と書けることを用いると，
    \begin{align}
      \sum_{n}\qty[(n + \alpha)(n + \alpha - 1) + (n + \alpha) - \nu^2]u_nz^{n + \alpha - 2} &= -\sum_{n}u_nz^{n + \alpha} \\ 
      \Leftrightarrow \sum_{n}\qty[(n + \alpha)(n + \alpha - 1) + (n + \alpha) - \nu^2]u_nz^{n + \alpha - 2} &= -\sum_{n}u_{n - 2}z^{n + \alpha -2} 
    \end{align}
    となる．係数を比較すると，
    \begin{align}
      \qty[(n + \alpha)^2 - \nu^2]u_n = -u_{n - 2}\label{bessel-func-recurr}
    \end{align}
    となる．
    \refe{bessel-func-recurr}に$n = 1$を代入すると$u_{-1} = 0$より，
    \begin{align}
      0 = u_1 = u_3 = \cdots
    \end{align}
    を得る．
    また，$n = 0$を代入すると$u_0 \neq 0$より，
    \begin{align}
      (0 + \alpha)^2 - \nu^2 &= 0 \\ 
      \Rightarrow \alpha = \pm \mu
    \end{align}
    となる．
    $\alpha = \nu$を採用して\refe{bessel-func-recurr}を用いると，
    \begin{align}
      u_{n + 2} = -\frac{1}{(n + \nu + 2)^2 - \nu^2}u_n 
    \end{align}
    を得る．
    $\nu \in \mathbb{Z}$のときは，
    \begin{align}
      J_{\nu}(z) &= u_0z^{\nu}\qty(1 - \frac{1}{2(2\nu + 2)}z^2 + \frac{1}{2\cdot 4(2\nu + 2)(2\nu + 4)}z^4 + \cdots + (-1)^n\frac{(2\nu)!!}{(2n)!!(2\nu + 2n)!!}z^{2n} + \cdots) \\ 
      &= u_0z^{\nu}\sum_{n = 0}(-1)^n\frac{2^{\nu}\nu!}{2^nn!(\nu + n)!2^{\nu + n}} z^{2n} \\ 
      &= u_0\nu!2^{\nu}\sum_{n = 0}\frac{1}{n!(\nu + n)!}\qty(\frac{z}{2})^{\nu + 2n}
    \end{align}
    となる．$\nu \notin \mathbb{Z}$のときも，表せるようにガンマ関数をもちいて一般化する．
    \refe{bessel-differential-eq}の形より，明らかに定数倍が許容されるから，
    \begin{align}
      u_0 = \frac{1}{2^{\nu}\Gamma(\nu + 1)}
    \end{align}
    となるように$u_0$を定めておくと，
    \begin{align}
      J_{\nu}(z) &= \sum_{n}\frac{(-1)^n}{n!\Gamma(\nu + k + 1)}\qty(\frac{z}{2})^{\nu + 2n}
    \end{align}
    を得る．$J_{\nu}(z)$をBessel関数という．
    また，$\alpha = -\nu$を採用したときは，
    \begin{align}
      J_{-\nu}(z) &= \sum_{n}\frac{(-1)^n}{n!\Gamma(-\nu + k + 1)}\qty(\frac{z}{2})^{-\nu + 2n}
    \end{align}
    となり，これは$\nu \notin \mathbb{Z}$のときに$J_{\nu}(z)$と独立な解となることが示せる．
    $J_{-\nu}(z)$をNeumann関数という．
    \par
    得られたBessel関数とNeumann関数を用いて，
    \begin{align}
      j_n(z) &\coloneqq \sqrt{\frac{\pi}{2z}}J_{n + \frac{1}{2}}(z) \\ 
      y_n(z) &\coloneqq (-1)^{n + 1}\sqrt{\frac{\pi}{2z}}J_{-n-\frac{1}{2}}(z) \\ 
    \end{align}
    を定義する．$j_n(z)$と$y_n(z)$はそれぞれ，球Bessel関数，球Neumann関数という．
    さらに，球Bessel関数と球Neumann関数を用いて，
    \begin{align}
      h_n^{(1)}(z) &\coloneqq j_n(z) + \i y_n(z) \\ 
      h_n^{(2)}(z) &\coloneqq j_n(z) - \i y_n(z)
    \end{align}
    を定義する．$h_n^{(1)}$と$h_n^{(2)}$はそれぞれ，第1種Hankel関数，第2種Hankel関数という．
  \subsection{Helmholtz方程式}
    Helmholtz方程式，
    \begin{align}
      \qty(\laplacian + \kappa^2)u(r, \theta, \phi) = 0\label{helmholtz-eq-wave}
    \end{align}
    を考える．
    $u(\bm{r})$を正規直交基底である$\e^{\i m\phi}$で展開すると，
    \begin{align}
      A_m(r, \theta) &\coloneqq \int_{0}^{2\pi}u(r, \theta, \phi)\e^{-\i m\phi}\dd{\phi} \\ 
      u(r, \theta, \phi) &= \sum_{m}A_m(r, \theta)\e^{\i m\phi}\label{u-expand-with-a} \\ 
    \end{align}
    となる．
    また，極座標ラプラシアンは，
    \begin{align}
      \laplacian = \frac{1}{r^2}\pdv{r}\qty(r^2\pdv{r}) + \frac{1}{r^2\sin\theta}\pdv{\theta}\qty(\sin\theta\pdv{\theta}) + \frac{1}{r^2\sin^2\theta}\pdv[2]{\phi} \label{polar-laplacian}
    \end{align}
    である．
    \refe{helmholtz-eq-wave}に\refe{polar-laplacian}を用いて，\refe{u-expand-with-a}を代入すると，
    \begin{align}
      &\frac{1}{r^2}\pdv{r}\qty(r^2\pdv{r}\sum_{m}A_m(r, \theta)\e^{\i m\phi}) + \frac{1}{r^2\sin\theta}\pdv{\theta}\qty(\sin\theta\pdv{\theta}\sum_{m}A_m(r, \theta)\e^{\i m\phi}) + \frac{1}{r^2\sin^2\theta}\pdv[2]{\phi}\sum_{m}A_m(r, \theta)\e^{\i m\phi} \notag \\ 
      &+ \kappa^2\sum_{m}A_m(r, \theta)\e^{\i m\phi} = 0 \\ 
      \Leftrightarrow &\sum_{m}\e^{\i m\phi}\qty[
        \frac{1}{r^2}\pdv{r}\qty(r^2\pdv{r}A_m(r, \theta)) + 
        \frac{1}{r^2\sin\theta}\pdv{\theta}\qty(\sin\theta\pdv{\theta}A_m(r, \theta)) 
        - \frac{m^2}{r^2\sin^2\theta}A_m(r, \theta) + \kappa^2A_m(r, \theta)
      ] = 0 \label{am-before-int}
    \end{align}
    となる．
    \refe{am-before-int}に左から$\int_{0}^{2\pi}\dd{\phi}\e^{-\i m\phi}$をかける，すなわち，$\e^{\i m\phi}$に射影すると，
    \begin{align}
      \frac{1}{r^2}\pdv{r}\qty(r^2\pdv{r}A_m(r, \theta)) + \frac{1}{r^2\sin\theta}\pdv{\theta}\qty(\sin\theta\pdv{\theta}A_m(r, \theta)) - \frac{m^2}{r^2\sin^2\theta}A_m(r, \theta) + \kappa^2A_m(r, \theta) = 0 \label{a-condition}
    \end{align}
    を得る．
    Legendre陪多項式を定義するときに用いた$\mathcal{L}_m$において，$x = \cos\theta$とすると，
    \begin{align}
      \mathcal{L}_m &= \dv{x}\qty{(1 - \cos^2\theta)\dv{x}} - \frac{m^2}{1 - \cos^2\theta} \\ 
      &= \dv{x}\qty{\sin^2\theta\dv{x}} - \frac{m^2}{\sin^2\theta} \\ 
      &= \dv{\theta}{x}\dv{x}\qty{\sin^2\theta\dv{\theta}{x}\dv{x}} - \frac{m^2}{\sin^2\theta} \\ 
      &= -\frac{1}{\sin\theta}\dv{x}\qty{\sin^2\theta\qty(-\frac{1}{\sin\theta})} - \frac{m^2}{\sin^2\theta} \\ 
      &= \frac{1}{\sin\theta}\dv{x}\qty{\sin\theta\dv{\theta}} - \frac{m^2}{\sin^2\theta}
    \end{align}
    であるから，\refe{a-condition}は，
    \begin{align}
      \frac{1}{r^2}\pdv{r}\qty(r^2\pdv{r}A_m(r, \theta)) + \frac{1}{r^2}\mathcal{L}_{\abs{m}}A_m(r, \theta) + \kappa^2 A_m(r, \theta) = 0 \label{a-condition-use-lm}
    \end{align}
    となる．
    \par
    続いて，$A_m(r, \theta)$を直交基底であるLegendre陪多項式を用いて，
    \begin{align}
      A_m(r, \theta) = \sum_{n = \abs{m}}^{\infty}B_{nm}(r)P_{n}^{\abs{m}}(\cos\theta) \label{am-expand-with-pnm}
    \end{align}
    と展開する．
    \refe{a-condition-use-lm}の両辺に$P_{n}^{\abs{m}}(\cos\theta)$との内積をとる．
    \refe{pnm-norm}より，
    \begin{align}
      \inner{P_n^{m}(\cos\theta)}{P_n^{m}(\cos\theta)} = \delta_{n}^{n'}\frac{2}{2n + 1}\frac{(n + m)!}{(n - m)!}
    \end{align}
    であることと，\refe{assos-legendre-diff-eq}より，
    \begin{align}
      \mathcal{L}_mP_n^m(x) + n(n + 1)P_n^m(x) = 0
    \end{align}
    であることを用いれば，
    \begin{align}
      & \inner{P_n^{m}(\cos\theta)}{\frac{1}{r^2}\pdv{r}\qty(r^2\pdv{r}A_m(r, \theta)) + \frac{1}{r^2}\mathcal{L}_{\abs{m}}A_m(r, \theta) + \kappa^2 A_m(r, \theta)} = 0 \\ 
      \Leftrightarrow & \frac{1}{r^2}\inner{P_n^{m}(\cos\theta)}{\pdv{r}\qty(r^2\pdv{r}A_m(r, \theta))} + \frac{1}{r^2}\inner{P_n^{m}(\cos\theta)}{\mathcal{L}_{\abs{m}}A_m(r, \theta)} + \inner{P_n^{m}(\cos\theta)}{\kappa^2 A_m(r, \theta)} = 0 \\ 
      \Leftrightarrow & \frac{1}{r^2}\qty{\dv{r}r^2\qty{\dv{r}B_{nm}(r)}}\inner{P_n^{m}(\cos\theta)}{P_n^{m}(\cos\theta)} + \frac{1}{r^2}B_{nm}(r)\inner{P_n^{m}(\cos\theta)}{\mathcal{L}_{\abs{m}}P_n^{m}(\cos\theta)} \notag \\ 
      &+ \kappa^2B_{nm}(r)\inner{P_n^{m}(\cos\theta)}{P_n^{m}(\cos\theta)} = 0 \\ 
      \Leftrightarrow&  \qty(\frac{1}{r^2}\qty{\dv{r}r^2\qty{\dv{r}B_{nm}(r)}} - \frac{1}{r^2}n(n + 1)B_{nm}(r) + \kappa^2 B_{nm}(r))\inner{P_n^{m}(\cos\theta)}{P_n^{m}(\cos\theta)} = 0 \\ 
      \Rightarrow& \frac{1}{r^2}\qty{\dv{r}r^2\qty{\dv{r}B_{nm}(r)}} - \frac{1}{r^2}n(n + 1)B_{nm}(r) + \kappa^2 B_{nm}(r) = 0\label{bnm-condition}
    \end{align}
    を得る．
    \par
    \refe{bnm-condition}において，
    \begin{align}
      \rho \coloneqq \kappa r\label{rho-def}
    \end{align}
    と変数変換して，
    \begin{align}
      C_{nm}(\rho) \coloneq \sqrt{\rho}B_{nm}(r)\label{cnm-def}
    \end{align}
    と定義して代入すれば，
    \begin{align}
      \qty{\dv[2]{\rho} + \frac{1}{\rho}\dv{\rho} + 1 - \frac{1}{\rho^2}\qty(n + \frac{1}{2})^2}C_{nm}(\rho) = 0\label{bessel-form-rho}
    \end{align}\
    を得る．
    \refe{bessel-form-rho}と\refe{bessel-differential-eq-sub}である，
    \begin{align}
      \mathcal{L}_{\nu}J_{\nu}(z) &= - J_{\nu}(z)  \\ 
      \Leftrightarrow \qty(\dv[2]{z} + \frac{1}{z}\dv{z} + 1 - \frac{\nu^2}{z^2}) &= 0 \label{bessel-differential-eq-sub-2}
    \end{align}
    を比べると，\refe{bessel-differential-eq-sub-2}において，
    \begin{align}
      z &\to x \\ 
      \nu &= n + \frac{1}{2}
    \end{align}
    としたものであると分かる．
    つまり，$C_{nm}(\rho)$はBessel関数$J_{n + \frac{1}{2}}(\rho)$とNeumann関数$J_{-n - \frac{1}{2}}(\rho)$の線型結合で書かれるので，
    係数をそれぞれ$a'_{nm}$，$b'_{nm}$として，
    \begin{align}
      C_{nm}(\rho) = a'_{nm}J_{n + \frac{1}{2}}(\rho) + b'_{nm}J_{-n - \frac{1}{2}}(\rho)\label{expand-cnm-ans-tmp}
    \end{align}
    と書く．
    今後の変形のために，球Bessel関数$j_n(\rho)$と球Neumann関数$y_n(\rho)$を用いて\refe{expand-cnm-ans-tmp}を書き直せば，
    \begin{align}
      C_{nm}(\rho) &= a'_{nm}J_{n + \frac{1}{2}}(\rho) + b'_{nm}J_{-n - \frac{1}{2}}(\rho) \\ 
      &= \sqrt{\frac{2\rho}{\pi}}a'_{nm}\pi^{-1}j_n(\rho) + \sqrt{\frac{2\rho}{\pi}}b'_{mn}y_n(\rho) \\ 
      &= \sqrt{\rho}\qty[a_{nm}j_n(\rho) + b_{mn}y_n(\rho)] \label{expand-cnm-ans}
    \end{align}
    である．ただし，
    \begin{align}
      a_{nm} &\coloneq a'_{nm}\sqrt{\frac{2}{\pi}} \\ 
      b_{nm} &\coloneq b'_{nm}\sqrt{\frac{2}{\pi}}
    \end{align}
    と定めた．
    \par
    さて，今Helmholtz方程式の解の基底展開を行っているのであった．
    今までに定義した展開を全てまとめると，
    \begin{align}
      \text{\refe{u-expand-with-a}より}&\ u(\bm{r}) = \sum_{m = -\infty}^{\infty} A_m(r, \theta) \e^{\i m\phi} \\ 
      \text{\refe{am-expand-with-pnm}より}&\ A_m(r, \theta) = \sum_{n = \abs{m}}^{\infty}B_{nm}(r)P_{n}^{\abs{m}}(\cos\theta) \\ 
      \text{\refe{rho-def}より}&\ \rho = \kappa r \\ 
      \text{\refe{cnm-def}より}&\ B_{nm}(r) = \frac{1}{\sqrt{\rho}}C_{nm}(\rho) \\ 
      \text{\refe{expand-cnm-ans}より}&\ C_{nm}(\rho) = \sqrt{\rho}\qty[a_{nm}j_n(\rho) + b_{mn}y_n(\rho)]
    \end{align}
    となる．下から上に全て代入して整理すると，
    \begin{align}
      u(\bm{r}) = \sum_{m = -\infty}^{\infty}\sum_{n = \abs{m}}^{\infty}\qty[a_{nm}j_n(\kappa r) + b_{nm}y_n(\kappa r)] P_{n}^{\abs{m}}(\cos\theta)\e^{\i m\phi}
    \end{align}
    となる．
\end{document}